\chapter{The operating discipline}
\label{ch:operating-discipline}

Every theorem in the correspondence declares its $(\text{level},
\text{chart}, \text{ambient})$ coordinates. Promotion across coordinates
requires the named comparison arrow, constructed under the named
hypotheses. No claim is permitted to be promoted from one coordinate to
another by elision.

The vertical axis is the universal arrow with five levels:
primitive at $0$ ($\mathrm{CY}_d$-categories with cyclic $\Ainf$-data),
Stage-1 native at $1$
($\PhiFA_d(\cC) \in E_d\text{-}\HolFA(X)$),
Stage-2 chart shadow at $2$
($A_X = \SpCh_{\Sigma_{d-1}, C} \!\circ \PhiFA_d(\cC)$),
centre / quantum vertex group at $3$
($Z^{\mathrm{der}}_{\mathrm{ch}}(A_X) \simeq \mathrm{ChirHoch}^\bullet(A_X, A_X)$,
the positive half $Y^+(X)$, the Drinfeld double $G(X) = D(Y^+(X))$),
scalar trace at $4$
($\kappa_{\mathrm{BKM}}(\Phi_N) = c_N(0)/2$).
Each transition is mediated by a named arrow:
Kontsevich--Tamarkin formality plus
Costello--Gwilliam--Li locality from $0$ to $1$, factorisation
homology along $\Sigma_{d-1}$ from $1$ to $2$, the bar/twisting
comparison $Z^{\mathrm{der}}_{\mathrm{ch}}(A) = \RHom(\Omega B(A), A)$
from $2$ to $3$, and the trace pairing from $3$ to $4$.

The horizontal axis is the chart datum, with internal product structure
\[
\text{Chart} \;=\; \text{equivariance stratum}
  \times_{\text{stratum}} (\Sigma_{d-1}, C)
  \times_{\text{stratum}} (\text{boundary vacuum } b)
  \times_{\text{stratum}} (\text{admissibility window}).
\]
The four equivariance strata are toric $T^d$, reduced
$\mathbb{C}^\times +$ $\mathrm{Aut}(X)$, orbifold inertia $I(X/G)$,
and lattice-polarised period domain.

The ambient axis is the depth context: ordinary chain complex,
weight-completed, pro, $J$-adic, HS-sewing, formal-local,
global-with-descent, derived $\infty$-categorical. The chain-level lane
and the $(\infty, 1)$-categorical lane are equally load-bearing; state
each theorem in the lane in which its proof actually works.

The discipline is detailed in
Appendix~\ref{app:three-axis-scope-discipline}; the open comparison
problems exposed by it are inscribed in
Chapter~\ref{ch:open-frontiers}. The statements throughout the
remainder of this part carry $(\text{level}, \text{chart},
\text{ambient})$ coordinates inline.
